\documentclass[10pt,a4paper]{amsart}
\usepackage[margin=19mm]{geometry}
\usepackage{amsmath,amssymb,mathtools}
\usepackage[scr=rsfs,cal=euler]{mathalfa}
\usepackage{xfrac}
\usepackage[unicode,hidelinks,bookmarks=false]{hyperref}
\newcommand{\ie}{i.e.\ }
\newcommand{\cf}{cf.\ }
\newcommand{\resp}{respectively\ }
\newcommand{\Subsection}{\subsection}
\providecommand{\nobreakdash}{\nobreak\hbox{-}}
\setlength{\emergencystretch}{2em}
\multlinegap0pt
\usepackage{mathtools}
\usepackage{bbm}
\usepackage{enumitem}
\usepackage{ulem}
\usepackage{mathrsfs}
\usepackage{pdflscape}
\usepackage{tikz,float}
\usepackage{pgfplots}
\pgfplotsset{compat=1.18}
\newtheorem{defi}{Definition}
\newtheorem{assumption}{Assumption}
\newtheorem{thm}{Theorem}
\newtheorem*{thm*}{Theorem}
\newtheorem{prop}{Proposition}
\newtheorem{lemma}{Lemma}
\newtheorem{remark}{Remark}
\newtheorem{coro}{Corollary}
\newtheorem{conj}{Conjecture}
\DeclareMathOperator{\diam}{diam}
\begin{document}
\title[Prescribed realisation of longest runs in continued fraction]{Prescribed realisation of longest runs in continued fractions}
\author{Ying Wai Lee}
\date{}
\maketitle
\noindent\emph{Source and licence.} Prescribed realisation of longest runs in continued fractions. Version arXiv:2606.00505v1 (CC BY 4.0). This material is distributed under the Creative Commons Attribution 4.0 International licence, http://creativecommons.org/licenses/by/4.0/, as recorded in the arXiv OAI licence metadata for this exact version and shown on the abstract page https://arxiv.org/abs/2606.00505v1. Full rights notice: licenses/arxiv-2606.00505-CC-BY-4.0.txt.\par\smallskip
\noindent\textit{Editorial scope.} Counted unit: the original \texttt{thm} environment \texttt{thm:main-sn-final} at source lines 143--158 together with its complete proof at lines 787--825; the delivered document keeps source lines 132--826 so that every referenced label and every citation resolves inside it. Native editable source is the paper's own concatenated TeX; the AMS wrapper is editorial formatting only. No publisher class, upstream script or unrelated artwork is redistributed.
\par\smallskip\noindent\textit{Reference note.} This article ships no compiled bibliography (biblatex); the entries delivered here are composed from the author's own BibTeX (.bib) fields, with no field invented and no entry dropped.
\section{Main results}

The following definition isolates, in functional notation, the exact hypothesis appearing in the above result of Wang--Wu~\cite[Theorem~1.3]{WangWu2011}.
\begin{defi}
Let $f:\mathbb{N}\to\mathbb{N}$ be a function. $f$ is said to be admissible, if $f$ is non-decreasing, unbounded, and
\begin{align*}
    \lim_{n\to\infty}\frac{f(n+f(n))}{f(n)}=1.
\end{align*}
\end{defi}

Theorem~\ref{thm:main-sn-final} is the main theorem of the present work.

\begin{thm}
\label{thm:main-sn-final}
For any $\lambda\in\mathbb{N}$ and admissible $f:\mathbb{N}\to\mathbb{N}$, 
\begin{align*}
    \dim{\mathfrak{E}_{\lambda,f}}=1,
\end{align*}
where
\begin{align*}
    \mathfrak{E}_{\lambda,f}
    &\coloneqq
    \left\{x\in[0,1)\setminus\mathbb{Q}:
    \lim_{n\to\infty}\frac{L_n(x,\lambda)}{f(n)}=1
    \text{ and }\right.\\
    &\qquad\qquad\left.\text{for any $n\in\mathbb{N}$ and $\mu\in\mathbb{N}\setminus\{\lambda\}$, }\frac{L_n(x,\lambda)}{L_n(x,\mu)}>1\right\}.
\end{align*}
\end{thm}


The first corollary records the fixed-symbol exact-growth consequence of Theorem~\ref{thm:main-sn-final}. Compared with the fixed-symbol exceptional-set theorem of Song--Zhou, it gives a refinement in which the prescribed growth scale is allowed to be any admissible function. Thus this corollary isolates the fixed-symbol part of the main theorem after the additional dominance condition is discarded.
\begin{coro}
\label{label:coro1}
For any $\lambda\in\mathbb{N}$ and admissible $f:\mathbb{N}\to\mathbb{N}$, 
\begin{align*}
    \dim{\left\{x\in[0,1)\setminus\mathbb{Q}:
    \lim_{n\to\infty}\frac{L_n(x,\lambda)}{f(n)}=1\right\}}=1.
\end{align*}
\end{coro}

The second corollary records the overall exact-growth consequence of Theorem~\ref{thm:main-sn-final}. Compared with the theorem of Wang--Wu for the overall longest-run function, it recovers the same full-dimensional exact-growth statement. Thus this corollary shows that the main theorem contains the known overall result as a special consequence, with the additional information that the maximising symbol may be prescribed.
\begin{coro}
\label{label:coro2}
For any admissible $f:\mathbb{N}\to\mathbb{N}$, 
\begin{align*}
    \dim{\left\{x\in[0,1)\setminus\mathbb{Q}:
    \lim_{n\to\infty}\frac{R_n(x)}{f(n)}=1\right\}}=1.
\end{align*}
\end{coro}

The third corollary records the prescribed-symbol dominance consequence of Theorem~\ref{thm:main-sn-final}. Unlike the previous fixed-symbol and overall exact-growth results, this statement has no direct analogue in the earlier literature. Thus this corollary gives a new full-dimensional phenomenon: the partial quotient value responsible for the overall longest run can be fixed in advance.
\begin{coro}
\label{label:coro3}
For any $\lambda\in\mathbb{N}$,
\begin{align*}
    \dim{\left\{x\in[0,1)\setminus\mathbb{Q}:
    \text{for any $n\in\mathbb{N}$ and $\mu\in\mathbb{N}\setminus\{\lambda\}$, }\frac{
    L_n(x,\lambda)}{L_n(x,\mu)}>1\right\}}=1.
\end{align*}
\end{coro}

\section{Basic lemmas}

Let $\mathbb{N}^*\coloneqq\{\emptyset\}\cup\bigcup_{n\in\mathbb{N}}\mathbb{N}^n$ be the finite word space. The standard convention is adopted that juxtaposition of two or more finite words denotes their concatenation. If a finite word is put inside continued-fraction brackets, the brackets denote the finite continued fraction whose entries are the letters of that word. The same convention is used when a finite word is followed by an infinite continued-fraction tail.

Define $\mathcal{I}(\emptyset)\coloneqq[0,1]$ and for any $n\in\mathbb{N}$ and $\omega=(\omega_1,\dots,\omega_n)\in\mathbb{N}^n$, $\mathcal{I}(\omega)$ to be the continued-fraction cylinder of $\omega$:
\begin{align*}
    \mathcal{I}(\omega)
    \coloneqq
    \bigcap_{i=1}^n\overline{\left\{x\in[0,1)\setminus\mathbb{Q}:
    a_i(x)=\omega_i\right\}}.
\end{align*}
Define $K:\mathbb{N}^*\to\mathbb{N}$ to be the continuant function; that is $K(\emptyset)\coloneqq1$ and for any $n\in\mathbb{N}$ and $\omega=(\omega_1,\dots,\omega_n)\in\mathbb{N}^n$, $K(\omega)$ is the denominator of the fraction $[\omega_1,\ldots,\omega_n]$.

\begin{lemma}
\label{lem:continuant-concatenation}
For any $\omega_1,\omega_2\in\mathbb{N}^*$, $K(\omega_1)K(\omega_2)\leq K(\omega_1\omega_2)\leq 2K(\omega_1)K(\omega_2)$.
\end{lemma}
\begin{proof}
The claim is immediate when one of the words is empty; it suffices to handle the non-trivial case. By the continuant identity, one obtains $K(\omega_1\omega_2)=K(\omega_1)K(\omega_2)+K({\omega_{1}}\partial)\,K(\partial\omega_2)$,
where ${\omega_1}\partial$ denotes $\omega_1$ with its last digit removed and $\partial\omega_2$ denotes $\omega_2$ with its first digit removed. The result follows from $0\leq K({\omega_1}\partial)\leq K(\omega_1)$ and $0\leq K(\partial\omega_2)\leq K(\omega_2)$.
\end{proof}

Define, for any $\lambda\in\mathbb{N}$, $\lambda^0\coloneqq\emptyset$ and for $n\in\mathbb{N}$, the long-run word $\lambda^n\coloneqq(\lambda,\lambda,\ldots,\lambda)\in\mathbb{N}^n$.

\begin{lemma}
\label{lem:continuant-growth}
For any $\omega\in\mathbb{N}^n$, $K(\omega)\geq \varphi^{n-1}$. For any $\lambda\in\mathbb{N}$ and $n\in\mathbb{N}\cup\{0\}$, $K(\lambda^n)\leq {\tau_\lambda}^n$.
\end{lemma}
\begin{proof}
For any $n\in\mathbb{N}$ and $\omega\in\mathbb{N}^n$, $K(\omega)\geq K(1^n)\geq \varphi^{n-1}$. One obtains from the standard continuant recurrence formula that for any $n\in\mathbb{N}$, $K(\lambda^{n+1})=\lambda\,K(\lambda^n)+K(\lambda^{n-1})$ and 
\begin{align*}
    K(\lambda^n)
    =\frac{{\tau_\lambda}^{n+1}-(-\tau_\lambda)^{-n-1}}{\sqrt{\lambda^2+4}}
    \leq{\tau_\lambda}^n.
\end{align*}
\end{proof}
Define $T_{\emptyset}\coloneqq\mathrm{id}_{[0,1]}$. Define, for any $n\in\mathbb{N}$ and $\omega=(\omega_1,\ldots,\omega_n)\in\mathbb{N}^n$, the prefix map $T_\omega:[0,1]\to[0,1]$ by, for any $\xi\in[0,1]$,
\begin{align*}
    T_\omega(\xi)\coloneqq \frac{p_n+\xi p_{n-1}}{q_n+\xi q_{n-1}},
\end{align*}
where $p_0\coloneqq0$, $q_0\coloneqq 1$ and for any $k\in\{1,\ldots,n\}$, $p_k$ and $q_k$ are coprime positive integers so that $p_k/q_k=[\omega_1,\ldots,\omega_k]$. 

\begin{lemma}
\label{lem:prefix-map-derivative}
For any $\omega\in\mathbb{N}^*$, one obtains that $T_\omega([0,1])=\mathcal{I}(\omega)$ and for any $\xi\in[0,1]$, 
\begin{align*}
    \frac{1}{4K(\omega)^2}
    \leq\left|T_\omega'(\xi)\right|
    \leq\frac{1}{K(\omega)^2};
\end{align*}
in particular, $T_\omega$ is bi-Lipschitz and preserves the Hausdorff dimension of subsets of $[0,1]$.
\end{lemma}
\begin{proof}
The claim is immediate when the word is empty; it suffices to handle the non-trivial case. Pick any $n\in\mathbb{N}$ and $\omega\in\mathbb{N}^n$. By differentiation and the standard formula $|p_{n-1}q_n-p_nq_{n-1}|=1$, one obtains that for any $\xi\in[0,1]$, $\left|T_\omega'(\xi)\right|=1/(q_n+\xi q_{n-1})^2$. Since $0\leq q_{n-1}\leq q_n=K(\omega)$, one obtains that for any $\xi\in[0,1]$, $K(\omega)\leq q_n+\xi q_{n-1}\leq 2K(\omega)$. Since $T_\omega'$ is bounded by positive constants on $[0,1]$, one obtains from applying the mean value theorem that $T_\omega$ is bi-Lipschitz.
\end{proof}

Define, for any $B\in\mathbb{N}$, 
\begin{align*}
    \mathcal{F}_B
    \coloneqq
    \bigcap_{n\in\mathbb{N}}
    \left\{x\in[0,1)\setminus\mathbb{Q}: a_n(x)\in\{1,\ldots,B\}\right\}.
\end{align*}
\begin{lemma}
\label{lem:bounded-prefix-dimension}
For any $B,n\in\mathbb{N}$ and $\nu\in\{1,\dots,B\}^n$,
\begin{align*}
    \dim{\mathcal{F}_B}
    =\dim{\left(\mathcal{I}(\nu)\cap\mathcal{F}_B\right)}.
\end{align*}
\end{lemma}
\begin{proof}
Pick any $n,B\in\mathbb{N}$ and $\nu\in\{1,\dots,B\}^n$. 
By the definition of the prefix map $T_\nu$, one obtains that for any $\xi\in{\mathcal{F}_B}$, $T_\nu(\xi)=[\nu_1,\ldots,\nu_n,a_1(\xi),a_2(\xi),\ldots]$; hence, $\mathcal{I}(\nu)\cap \mathcal{F}_B=T_\nu(\mathcal{F}_B)$. By Lemma~\ref{lem:prefix-map-derivative}, $T_\nu$ is bi-Lipschitz on $\mathcal{F}_B$ and preserves Hausdorff dimension.
\end{proof}

\begin{lemma}
\label{lem:holder-dimension-estimate}
Let $(X,d_X)$ and $(Y,d_Y)$ be metric spaces. Let $g:X\to Y$ be a surjective function. Suppose there exist $0<\alpha\leq 1$ and $C>0$ such that for any $x_1,x_2\in X$,
\begin{align*}
    d_Y(g(x_1),g(x_2))\leq C \left(d_X(x_1,x_2)\right)^\alpha.
\end{align*}
Then $\dim{Y}\leq\dim{X}/\alpha$.
\end{lemma}
\begin{proof}
Pick any $s>\dim{X}$ and $\varepsilon>0$. Since $\mathcal{H}^s(X)=0$, there exists an open cover $\mathcal{U}_{\varepsilon}$ of $X$ such that $\sum_{U\in\mathcal{U}_{\varepsilon}}(\diam{U})^s<\varepsilon$ and for any $U\in\mathcal{U}_{\varepsilon}$, $\diam{U}<\varepsilon$. Since $g$ is surjective, $(g(U))_{U\in\mathcal{U}_{\varepsilon}}$ covers $Y$. Then for any $U\in\mathcal{U}_{\varepsilon}$, $\diam{g(U)}\leq C\,(\diam{U})^\alpha$, and
\begin{align*}
    \sum_{U\in\mathcal{U}_{\varepsilon}}(\diam{g(U)})^{s/\alpha}
    \leq C^{s/\alpha}\sum_{U\in\mathcal{U}_{\varepsilon}} (\diam{U})^s
    < C^{s/\alpha}\varepsilon.
\end{align*}
Hence, for any $s>\dim{X}$, $\mathcal{H}^{s/\alpha}(Y)=0$ and $\dim{Y}\leq s/\alpha$.
\end{proof}

The next lemma is a reformulation of~\cite[Lemma~2.5]{WangWu2011} in the present functional notation.
\begin{lemma}
\label{lem:fixed-multiple-self-neglecting}
Let $f:\mathbb{N}\to\mathbb{N}$. 
Suppose $f$ is admissible. 
Then for any $C\in\mathbb{N}$,
\begin{align*}
    \lim_{n\to\infty}\frac{f(n+C f(n))}{f(n)}
    =1.
\end{align*}
\end{lemma}

\section{Construction of a Cantor-type set}

Let $\lambda\in\mathbb{N}$ and $f:\mathbb{N}\to\mathbb{N}$ be a function. Define $H_\lambda\coloneqq\lfloor{\log{\tau_{\lambda}}}/{\log{\varphi}}\rfloor+1\in\mathbb{N}$. Pick any $M\in\mathbb{N}$. Define $t_0\coloneqq 1$ and for any $k\in\mathbb{N}$,
\begin{align*}
    r_k\coloneqq f(t_{k-1}),\qquad
    m_k\coloneqq \left\lfloor\sqrt{r_k}\right\rfloor, \qquad
    g_k\coloneqq MH_\lambda r_k, \qquad
    p_k\coloneqq \left\lfloor \frac{g_k}{m_k}\right\rfloor, \qquad
    q_k\coloneqq g_k-p_km_k,
\end{align*}
and
\begin{align*}
    s_k\coloneqq f(t_{k-1}+g_k), \qquad 
    d_k\coloneqq 2p_k+2+s_k, \qquad 
    l_k\coloneqq g_k+d_k, \qquad 
    t_k\coloneqq t_{k-1}+l_k.
\end{align*}

Let $k\in\mathbb{N}$. The quantity $r_k$ is the target scale at the beginning of stage $k$, while $s_k$ is the target scale associated with the end of the retained free length $g_k$. The length $g_k$ specifies the number of free digits retained at stage $k$. The auxiliary length $m_k$ tends to infinity while satisfying $\lim_{k\to\infty}m_k/r_k=0$; it will be used to subdivide the free digits into shorter pieces, so that accidental runs remain negligible compared with the target scale. The integers $p_k$ and $q_k$ are the quotient and remainder obtained when the free length $g_k$ is divided into pieces of length $m_k$. Finally, $d_k$, $l_k$, and $t_k$ record respectively the non-free length to be added at stage $k$, the total length of stage $k$, and the starting position of stage $k+1$.

\begin{prop}
\label{prop:scale-inputs-final}
Let $\lambda\in\mathbb{N}$ and $f:\mathbb{N}\to\mathbb{N}$ be a function. 
Suppose $f$ is admissible. Then for any $M\in\mathbb{N}$,
\begin{align}
\label{eq:scale-inputs-final}
    \lim_{k\to\infty}\frac{s_k}{r_k}
    =\lim_{k\to\infty}\frac{s_{k-1}}{r_k}
    =1,
    \qquad
    \lim_{k\to\infty}\frac{p_k}{r_k}=0,
\end{align}
and there exists $K_{\lambda,f,M}\in\mathbb{N}$ such that for any $k\in\mathbb{N}$ if $k\geq K_{\lambda,f,M}$ then
\begin{align}
\label{eq:ell-eventual-bound-final}
    l_k\leq (MH_\lambda+2)\,r_k.
\end{align}
\end{prop}
\begin{proof}
Since $f$ is unbounded and non-decreasing, one obtains that $(t_k)_{k\in\mathbb{N}}$ is strictly increasing and $\lim_{k\to\infty}r_k=\lim_{k\to\infty}f(t_{k-1})=+\infty$. By Lemma~\ref{lem:fixed-multiple-self-neglecting}, one obtains
\begin{align*}
    \lim_{k\to\infty}\frac{f(t_{k-1}+MH_\lambda f(t_{k-1}))}{f(t_{k-1})}=1;
\end{align*}
hence, $\lim_{k\to\infty}{s_k}/{r_k}=1$. Since $\lim_{k\to\infty}r_k=+\infty$, one obtains $\lim_{k\to\infty}m_k=+\infty$.

Since for any $k\in\mathbb{N}$, $p_km_k\leq g_k=MH_\lambda r_k$, one obtains $0\leq {p_k}/{r_k}\leq {MH_\lambda}/{m_k}$; hence, $\lim_{k\to\infty}{p_k}/{r_k}=0$. Note that for any $k\in\mathbb{N}$,
\begin{align*}
    \frac{l_k}{r_k}
    =
    \frac{g_k+d_k}{r_k}
    =
    MH_\lambda+\frac{2p_k}{r_k}+\frac{2}{r_k}+\frac{s_k}{r_k};
\end{align*}
hence, $\lim_{k\to\infty}{l_k}/{r_k}=MH_\lambda+1$ and there exists $K_{\lambda,f,M}\in\mathbb{N}$ such that for any $k\in\mathbb{N}$ if $k\geq K_{\lambda,f,M}$ then~\eqref{eq:ell-eventual-bound-final} holds.

One obtains that for any $k\in\mathbb{N}$ if $k\geq K_{\lambda,f,M}+1$ then 
\begin{align*}
    t_{k-1}
    =t_{k-2}+l_{k-1}
    \leq t_{k-2}+(MH_\lambda+2)\,r_{k-1}
    =t_{k-2}+(MH_\lambda+2)\,f(t_{k-2}),
\end{align*}
and
\begin{align*}
    1\leq \frac{f(t_{k-1})}{f(t_{k-2})}\le
    \frac{f(t_{k-2}+(MH_\lambda+2)\,f(t_{k-2}))}{f(t_{k-2})}.
\end{align*}
By Lemma~\ref{lem:fixed-multiple-self-neglecting}, one obtains $\lim_{k\to\infty}{f(t_{k-1})}/{f(t_{k-2})}=1$; hence, $\lim_{k\to\infty}{r_{k-1}}/{r_k}=1$. Thus,
\begin{align*}
    \lim_{k\to\infty}\frac{s_{k-1}}{r_k}
    =\lim_{k\to\infty}\frac{s_{k-1}}{r_{k-1}}\lim_{k\to\infty}\frac{r_{k-1}}{r_k}
    =1.
\end{align*}
\end{proof}

Let $\lambda\in\mathbb{N}$ and $f:\mathbb{N}\to\mathbb{N}$ be a function. Pick any $M,B\in\mathbb{N}$. Suppose $B\geq\lambda+2$.
Define, for any $k\in\mathbb N$ and $U_k\in\{1,\dots,B\}^{g_k}$, the expanded block $U^+_k\in\{1,\dots,B\}^{l_k}$ by:
\begin{align*}
    U^+_k
    \coloneqq
    U_{k,1}\sigma_1\sigma_2U_{k,2}\sigma_1\sigma_2\cdots U_{k,p_k}\sigma_1\sigma_2U_{k,p_k+1}\sigma_1\lambda^{s_k}\sigma_2,
\end{align*}
where $\sigma_1\coloneqq \lambda+1$ and $\sigma_2\coloneqq \lambda+2$ are the separator digits, $U_{k,1},\dots,U_{k,p_k}\in \{1,\dots,B\}^{m_k}$ and $U_{k,p_k+1}\in \{1,\dots,B\}^{q_k}$ are the unique words satisfying 
\begin{align*}
    U_k=U_{k,1}U_{k,2}\cdots U_{k,p_k}U_{k,p_k+1},
\end{align*}
with the convention that $\{1,\ldots,B\}^0\coloneqq \{\emptyset\}$. Define
\begin{align*}
    \mathcal{K}_{\lambda,f,M,B}
    \coloneqq\left\{\left[U^+_1U^+_2U^+_3\cdots\right]:
    U_k\in\{1,\dots,B\}^{g_k}\text{ for all $k\in\mathbb{N}$.}\right\}
    \subset\mathcal{F}_B.
\end{align*}
Define $\zeta_{\lambda,f,M,B}:\mathcal{F}_B\to\mathcal{K}_{\lambda,f,M,B}$ by for any $y\in \mathcal{F}_B$,
\begin{align*}
    \zeta_{\lambda,f,M,B}(y)
    \coloneqq\left[U^+_{y,1}U^+_{y,2}U^+_{y,3}\cdots\right],
\end{align*}
where the partial quotients of $y$ are partitioned into consecutive words $U_{y,1},U_{y,2},U_{y,3},\ldots$ so that $y=[U_{y,1}U_{y,2}U_{y,3}\cdots]$ and for any $k\in\mathbb{N}$, $U_{y,k}\in\{1,\dots,B\}^{g_k}$. 

Define
\begin{align*}
    \mathfrak{F}_{\lambda,f}
    &\coloneqq
    \left\{x\in[0,1)\setminus\mathbb{Q}:
    \lim_{n\to\infty}\frac{L_n(x,\lambda)}{f(n)}=1
    \text{ and }\right.\\
    &\qquad\qquad\left.\text{for all sufficiently large $n\in\mathbb{N}$ and for any $\mu\in\mathbb{N}\setminus\{\lambda\}$, }\frac{L_n(x,\lambda)}{L_n(x,\mu)}>1\right\}.
\end{align*}

\begin{prop}
\label{prop:run-control-final}
Let $\lambda\in\mathbb{N}$ and $f:\mathbb{N}\to\mathbb{N}$ be a function. Suppose $f$ is admissible. Then for any $M,B\in\mathbb{N}$, if $B\geq \lambda+2$ then
\begin{align*}
    \mathcal{K}_{\lambda,f,M,B}\subset \mathfrak{F}_{\lambda,f}.
\end{align*}
\end{prop}
\begin{proof}
Pick any $x\in \mathcal{K}_{\lambda,f,M,B}$. Since $\mathcal{K}_{\lambda,f,M,B}=\zeta_{\lambda,f,M,B}(\mathcal{F}_B)$, there exists $y\in \mathcal{F}_B$ such that $x=\zeta_{\lambda,f,M,B}(y)$. For any $k\in\mathbb{N}$, the expanded block $U^+_{y,k}$ has the following form
\begin{align*}
    U^+_{y,k}
    =U_{y,k,1}\sigma_1\sigma_2U_{y,k,2}\sigma_1\sigma_2\cdots U_{y,k,p_k}\sigma_1\sigma_2U_{y,k,p_k+1}\sigma_1\lambda^{s_k}\sigma_2,
\end{align*}
where $U_{y,k,1},\dots,U_{y,k,p_k}\in \{1,\dots,B\}^{m_k}$ and $U_{y,k,p_k+1}\in \{1,\dots,B\}^{q_k}$ are the unique words satisfying $U_{y,k}=U_{y,k,1}U_{y,k,2}\cdots U_{y,k,p_k}U_{y,k,p_k+1}$. Pick any $k\in\mathbb{N}$. Every consecutive run in the first $t_k-1$ digits of $x$ which is not contained in one of the prescribed words $\lambda^{s_1},\ldots,\lambda^{s_k}$ lies either inside one of the words $U_{y,j,i}$ for some $j\in\{1,\ldots,k\}$ and $i\in\{1,\ldots,p_j+1\}$, or across a boundary of one of the forms
\begin{align*}
    U_{y,j,i}\sigma_1,\qquad
    \sigma_2U_{y,j,i+1},\qquad
    U_{y,j,p_j+1}\sigma_1,\qquad
    \sigma_2U_{y,\min\{j+1,k\},1},
\end{align*}
for some $j\in\{1,\ldots,k\}$ and $i\in\{1,\ldots,p_j\}$. Since $f$ is non-decreasing, the sequence $(m_k)_{k\in\mathbb{N}}$ is non-decreasing. Hence, each such consecutive run has length at most $m_k+1$. 

Since $\lim_{k\to\infty}r_k=\lim_{k\to\infty}m_k=+\infty$ and for any $k\in\mathbb{N}$, 
\begin{align*}
    0
    \leq \frac{m_k+1}{r_k}=\frac{m_k}{r_k}+\frac{1}{r_k}
    \leq \frac{1}{m_k}+\frac{1}{r_k},
\end{align*}
one obtains $\lim_{k\to\infty}{(m_k+1)}/{r_k}=0$. By~\eqref{eq:scale-inputs-final}, one obtains $\lim_{k\to\infty}{(m_k+1)}/{s_{k-1}}=0$ and that there exists $k_0\in\mathbb{N}$ such that $k_0\geq2$ and for any $k\in\mathbb{N}$, if $k\geq k_0$ then $m_k+1<s_{k-1}$. 

One obtains that for any $k\in\mathbb{N}$, $t_k=t_{k-1}+l_k>t_{k-1}+g_k$; hence, the sequence $(t_{k-1}+g_k)_{k\in\mathbb{N}}$ is strictly increasing. Since $f$ is non-decreasing, the sequence $(s_k)_{k\in\mathbb{N}}$ is unbounded and non-decreasing. One obtains that there exists $k_1\in\mathbb{N}$ such that $k_1\geq k_0+1$ and for any $k\in\mathbb{N}$, if $k\geq k_1$ then
\begin{align*}
    C_0\coloneqq\max{\left\{s_{k_0-1},m_{k_0-1}+1\right\}}<s_{k-1}.
\end{align*}

Pick any $k,n\in\mathbb{N}$. Suppose $k\geq k_1$ and $t_{k-1}\leq n< t_k$. Since the first $n$ partial quotients of $x$ contain the full block $\lambda^{s_{k-1}}$, one obtains $L_n(x,\lambda)\geq s_{k-1}$. On the other hand, every $\lambda$-run outside the prescribed blocks has length at most
\begin{align*}
    \max{\left\{C_0,m_k+1\right\}}<s_{k-1}\leq s_k.
\end{align*}
Every prescribed block from an earlier stage has length at most $s_{k-1}$, and the part of the current prescribed block already contained in the first $n$ partial quotients of $x$ has length at most $s_k$. Therefore, for any $k,n\in\mathbb{N}$, if $k\geq k_1$ and $t_{k-1}\leq n< t_k$ then
\begin{align}
    \label{eq:Ln-sandwich-final}
    s_{k-1}
    \leq L_n(x,\lambda)
    \leq s_k.
\end{align}

By Proposition~\ref{prop:scale-inputs-final}, there exists $K_{\lambda,f,M}\in\mathbb{N}$ such that for any $k\in\mathbb{N}$, if $k\geq K_{\lambda,f,M}$ then 
\begin{align*}
    t_k=t_{k-1}+l_k
    \leq t_{k-1}+(MH_\lambda+2)\,r_k
    =t_{k-1}+(MH_\lambda+2)\, f(t_{k-1}).
\end{align*}
One obtains that for any $k,n\in\mathbb{N}$, if $k\geq K_{\lambda,f,M}$ and $t_{k-1}\leq n< t_k$ then
\begin{align*}
    t_{k-1}\leq n< t_k\leq t_{k-1}+(MH_\lambda+2)\,f(t_{k-1});
\end{align*}
hence, one obtains from $f$ being non-decreasing that
\begin{align*}
    f(t_{k-1})
    \leq f(n)
    \leq f(t_{k-1}+(MH_\lambda+2)\,f(t_{k-1})).
\end{align*}
One obtains from Lemma~\ref{lem:fixed-multiple-self-neglecting} that
\begin{align*}
    \lim_{k\to\infty}\frac{f(t_{k-1}+(MH_\lambda+2)\,f(t_{k-1}))}{f(t_{k-1})}
    =1.
\end{align*}
By the definition of $(r_k)_{k\in\mathbb{N}}$ and the squeeze theorem, one obtains
\begin{align*}
    \lim_{k\to\infty}\sup_{n\in\{t_{k-1},\ldots,t_k-1\}}
    \left|\frac{f(n)}{r_k}-1\right|
    =0.
\end{align*}
Therefore, one obtains from combining~\eqref{eq:scale-inputs-final} and~\eqref{eq:Ln-sandwich-final} that
\begin{align*}
    \lim_{n\to\infty}\frac{L_n(x,\lambda)}{f(n)}=1.
\end{align*}

One remains to prove eventual dominance. Pick any $k,n\in\mathbb{N}$ and $\mu\in\mathbb{N}\setminus\{\lambda\}$. Suppose $k\geq k_1$ and $t_{k-1}\leq n< t_k$. Every $\mu$-run in the first $n$ partial quotients of $x$ is necessarily not one of the prescribed $\lambda$-blocks. Since $\mu\neq\lambda$, no $\mu$-run can be contained in a prescribed block $\lambda^{s_j}$. Hence, its length is bounded by the maximal length of a non-prescribed run, and one obtains
\begin{align*}
    L_n(x,\mu)\leq\max{\{C_0,m_k+1\}}<s_{k-1}\leq L_n(x,\lambda).
\end{align*}
Therefore, one obtains that for any $n\in\mathbb N$ and $\mu\in\mathbb{N}\setminus\{\lambda\}$, if $n\geq N_{\lambda,f,M}\coloneqq t_{k_1-1}$, then
\begin{align*}
    L_n(x,\lambda)>L_n(x,\mu).
\end{align*}
\end{proof}

\section{Continuant distortion}

Let $\lambda\in\mathbb{N}$. Define $C_\lambda\coloneqq8(\lambda+1)(\lambda+2)$. Let $M\in\mathbb{N}$ and $0<\alpha<M/(M+1)$. Note that $(1-\alpha)M>\alpha$ and $H_\lambda\log{\varphi}>\log{\tau_\lambda}$. Define
\begin{align*}
    \eta_{\lambda,M,\alpha}
    \coloneqq\frac{(1-\alpha)MH_\lambda\log\varphi-\alpha\log \tau_\lambda}{4}>0.
\end{align*}

\begin{lemma}
\label{lem:full-stage-comparison}
Let $\lambda\in\mathbb{N}$ and $f:\mathbb{N}\to\mathbb{N}$ be a function. Suppose $f$ is admissible. Then for any $M,B\in\mathbb{N}$ and $0<\alpha<M/(M+1)$, there exists $k_{\lambda,f,M,\alpha,1}\in\mathbb{N}$ such that for any $k\in\mathbb{N}$, if $k\geq k_{\lambda,f,M,\alpha,1}$ then for any $U_k\in\{1,\ldots,B\}^{g_k}$,
\begin{align*}
    K(U^+_k)^\alpha
    \leq 2^{-2\alpha}K(U_k).
\end{align*}
\end{lemma}
\begin{proof}
Pick any $k\in\mathbb{N}$. By Lemma~\ref{lem:continuant-concatenation}, one obtains
\begin{align*}
    K(U^+_k)
    \leq
    {C_\lambda}^{p_k+1}K(\lambda^{s_k})K(U_k).
\end{align*}
By Lemma~\ref{lem:continuant-growth}, one obtains
\begin{align*}
    \log{\frac{K(U^+_k)^\alpha}{K(U_k)}}
    \leq
    \alpha(p_k+1)\log C_\lambda
    +\alpha s_k\log\tau_\lambda
    -(1-\alpha)(g_k-1)\log\varphi.
\end{align*}
By the definition of $(g_k)_{k\in\mathbb{N}}$, one obtains
\begin{align*}
    \log{\frac{K(U^+_k)^\alpha}{K(U_k)}}
    \leq
    \alpha(p_k+1)\log C_\lambda+\alpha(s_k-r_k)\log\tau_\lambda
    -4\eta_{\lambda,M,\alpha} r_k+(1-\alpha)\log{\varphi}.
\end{align*}
By Proposition~\ref{prop:scale-inputs-final} and $\lim_{k\to\infty}r_k=+\infty$, one obtains that for all sufficiently large $k\in\mathbb{N}$, 
\begin{align*}
    \log{\frac{K(U^+_k)^\alpha}{K(U_k)}}\leq-2\eta_{\lambda,M,\alpha} r_k\leq -2\alpha\log 2.
\end{align*}
The result follows after exponentiating.
\end{proof}

\begin{lemma}
\label{lem:stage-prefix-comparison}
Let $\lambda\in\mathbb{N}$ and $f:\mathbb{N}\to\mathbb{N}$ be a function. Suppose $f$ is admissible. Then for any $M,B\in\mathbb{N}$ and $0<\alpha<M/(M+1)$, there exist $C_{\lambda,f,M,\alpha,2}>0$ and $k_{\lambda,f,M,\alpha,2}\in\mathbb N$ such that for any $k\in\mathbb{N}$ and $U_k\in\{1,\ldots,B\}^{g_k}$, if $k\geq k_{\lambda,f,M,\alpha,2}$ then for any prefix $W$ of $U^+_k$,
\begin{align*}
    K(W)^\alpha\leq C_{\lambda,f,M,\alpha,2}\,K(W^-),
\end{align*}
where $W^-$ is the underlying word obtained from $W$ by deleting all inserted symbols $\sigma_1,\sigma_2$, as specified by their positions in the construction, not all occurrences of the numerical values $\lambda+1$ or $\lambda+2$, and any terminal segment of the form $\sigma_1\lambda^u \sigma_2$ or $\sigma_1\lambda^u$ or $\sigma_1$ coming from the final inserted block.
\end{lemma}
\begin{proof}
Pick any $k\in\mathbb{N}$ and $U_k\in\{1,\ldots,B\}^{g_k}$. Let $W$ be a prefix of $U^+_k$.

Suppose $W$ does not enter the terminal block $\sigma_1\lambda^{s_k}\sigma_2$. Then $W$ consists of $j$ complete pieces $U_{k,i}\sigma_1\sigma_2$, followed by a final prefix lying in one of $U_{k,j+1}$, $U_{k,j+1}\sigma_1$, or $U_{k,j+1}\sigma_1\sigma_2$, with $0\leq j\leq p_k$. The corresponding retained word $W^-$ is a prefix of $U_k$, and $W^-$ is of length at least $jm_k$. By repeated use of Lemma~\ref{lem:continuant-concatenation}, each insertion of a block $\sigma_1\sigma_2$ between two retained subwords increases the continuant by at most a multiplicative factor $4K(\sigma_1\sigma_2)$, while appending a
terminal $\sigma_1$ or $\sigma_1\sigma_2$ contributes at most a factor $8K(\sigma_1)K(\sigma_2)$. Hence, for any $j\in\{0,\ldots,p_k\}$,
\begin{align*}
    K(W)\leq {C_\lambda}^{j+1}K(W^-).
\end{align*}
By Lemma~\ref{lem:continuant-growth}, one obtains that for any $j\in\{0,\ldots,p_k\}$,
\begin{align*}
    \log\frac{K(W)^\alpha}{K(W^-)}
    &\leq
    \alpha(j+1)\log C_\lambda
    -(1-\alpha)\log K(W^-) \\
    &\leq
    \alpha(j+1)\log C_\lambda-(1-\alpha)(j m_k-1)\log\varphi.
\end{align*}
Since $\lim_{k\to\infty}m_k=+\infty$, the coefficient of $j$ on the right-hand side is negative for all sufficiently large $k\in\mathbb{N}$. Hence, the right-hand side is bounded above by a constant depending only on $\lambda$, $M$ and $\alpha$.

Suppose $W$ enters the terminal block $\sigma_1\lambda^{s_k}\sigma_2$. In this case, there exist $u\in\{0,\ldots,s_k\}$ and $\theta\in\{\varnothing,\sigma_2\}$ such that
\begin{align*}
    W=
    U_{k,1}\sigma_1\sigma_2\cdots U_{k,p_k}\sigma_1\sigma_2U_{k,p_k+1}\sigma_1\lambda^{u}\theta;
\end{align*}
in particular $W^-=U_k$. By Lemma~\ref{lem:continuant-concatenation},
\begin{align*}
    K(W)\leq {C_\lambda}^{p_k+1}K(\lambda^u)K(U_k).
\end{align*}
By Lemma~\ref{lem:continuant-growth},
\begin{align*}
    \log\frac{K(W)^\alpha}{K(W^-)}
    =
    \log\frac{K(W)^\alpha}{K(U_k)}
    \leq
    \alpha(p_k+1)\log C_\lambda+\alpha(u+1)\log\tau_\lambda-(1-\alpha)(g_k-1)\log\varphi.
\end{align*}
By $u\leq s_k$, $g_k=MH_\lambda r_k$ and Proposition~\ref{prop:scale-inputs-final}, one obtains for all sufficiently large $k\in\mathbb{N}$, 
\begin{align*}
    \log\frac{K(W)^\alpha}{K(W^-)}
    \leq -2\eta_{\lambda,M,\alpha} r_k<0.
\end{align*}
\end{proof}

Let $\lambda,M,B\in\mathbb{N}$ and $f:\mathbb{N}\to\mathbb{N}$ be a function. Suppose $B\geq\lambda+2$. Define, for any $k\in\mathbb{N}$, $\mathscr{C}_k(\mathcal{K}_{\lambda,f,M,B})$ to be the family of stage-$k$ cylinders of $\mathcal{K}_{\lambda,f,M,B}$; that is,
\begin{align*}
    \mathscr{C}_k(\mathcal{K}_{\lambda,f,M,B})
    \coloneqq
    \left\{J=\mathcal{I}\!\left(U^+_1U^+_2\cdots U^+_k\right)\cap\mathcal{K}_{\lambda,f,M,B}:
    U_i\in\{1,\dots,B\}^{g_i}\text{ for all $i\in\{1,2,\ldots,k\}$.}\right\}.
\end{align*}
Define $\pi_{\lambda,f,M,B}:\mathcal{K}_{\lambda,f,M,B}\to \mathcal{F}_B$ to be the map obtained by deleting, in each stage, all inserted separators $\sigma_1\sigma_2$ together with the terminal block $\sigma_1\lambda^{s_k}\sigma_2$. Note that $\pi_{\lambda,f,M,B}$ is surjective and $\pi_{\lambda,f,M,B}\circ\zeta_{\lambda,f,M,B}=\mathrm{id}_{\mathcal{F}_B}$.

\begin{prop}
\label{prop:local-holder-final}
Let $\lambda\in\mathbb{N}$ and $f:\mathbb{N}\to\mathbb{N}$ be a function. Suppose $f$ is admissible. Then for any $M,B\in\mathbb{N}$ and $0<\alpha<M/(M+1)$, if $B\geq\lambda+2$ then there exist $k_{\lambda,f,M,\alpha}\in\mathbb{N}$ and $C_{\lambda,f,M,B,\alpha}>0$ such that for any
\begin{align*}
    J
    =\mathcal{I}\!\left(U^+_1U^+_2\cdots U^+_{k_{\lambda,f,M,\alpha}}\right)\cap\mathcal{K}_{\lambda,f,M,B}
    \in\mathscr{C}_{k_{\lambda,f,M,\alpha}}(\mathcal{K}_{\lambda,f,M,B}),
\end{align*}
where \(U_1,\ldots,U_{k_{\lambda,f,M,\alpha}}\) are the unique words determining \(J\), the restriction of $\pi_{\lambda,f,M,B}$ to $J$:
\begin{align*}
    \pi_{\lambda,f,M,B}|_J:J\to \mathcal{I}\!\left(U_1U_2\cdots U_{k_{\lambda,f,M,\alpha}}\right)\cap\mathcal{F}_B,
\end{align*}
is surjective and
for any $x,y\in J$,
\begin{align*}
    \left|\pi_{\lambda,f,M,B}(x)-\pi_{\lambda,f,M,B}(y)\right|
    \leq C_{\lambda,f,M,B,\alpha} \left|x-y\right|^\alpha.
\end{align*}
\end{prop}
\begin{proof}
The surjectivity of $\pi_{\lambda,f,M,B}|_J$ is immediate from the definition of $\pi$ and $\zeta_{\lambda,f,M,B}$.

Define $k_{\lambda,f,M,\alpha}\coloneqq \max{\{k_{\lambda,f,M,\alpha,1},k_{\lambda,f,M,\alpha,2}\}}$, where $k_{\lambda,f,M,\alpha,1}$ and $k_{\lambda,f,M,\alpha,2}$ are given by Lemmas~\ref{lem:full-stage-comparison} and~\ref{lem:stage-prefix-comparison} respectively.
Pick any $J\in\mathscr{C}_{k_{\lambda,f,M,\alpha}}(\mathcal{K}_{\lambda,f,M,B})$. There exist unique $U_{J,i}\in\{1,\ldots,B\}^{g_i}$ such that 
\begin{align*}
    J=\mathcal{I}(V^+_J)\cap\mathcal{K}_{\lambda,f,M,B},
\end{align*}
where $V^+_J\coloneqq U_{J,1}^+U_{J,2}^+\cdots U_{J,k_{\lambda,f,M,\alpha}}^+$. Define $V_J\coloneqq U_{J,1}U_{J,2}\cdots U_{J,k_{\lambda,f,M,\alpha}}$ and
\begin{align*}
    C_{\lambda,f,M,B,\alpha,0}
    \coloneqq 2^{2\alpha}C_{\lambda,f,M,\alpha,2}
    \max_{J\in\mathscr{C}_{k_{\lambda,f,M,\alpha}}(\mathcal{K}_{\lambda,f,M,B})}{\frac{K(V^+_J)^\alpha}{K(V_J)}}<+\infty,
\end{align*}
as the family $\mathscr{C}_{k_{\lambda,f,M,\alpha}}(\mathcal{K}_{\lambda,f,M,B})$ is finite. Define, for any $u\in\{1,\dots,B\}$, the interval $J_{B,u}$ by
\begin{align*}
    J_{B,u}\coloneqq
    \left[
    \frac{1}{u+1/\varphi},
    \frac{1}{u+1/\tau_B}
    \right].
\end{align*}
Note that for any $\xi\in\mathcal{F}_B$ and $u\in\{1,\ldots,B\}$, if $a_1(\xi)=u$ then $T(\xi)\in[1/\tau_B,1/\varphi]$ and $\xi\in J_{B,u}$. Thus, the family $\mathscr{J}_B\coloneqq(J_{B,u})_{u\in\{1,\ldots,B\}}$ consists of pairwise disjoint compact intervals, and one obtains:
\begin{align*}
    \Delta_B
    \coloneqq
    \min_{1\leq u_1<u_2\leq B}\operatorname{dist}{\left(J_{B,u_1},J_{B,u_2}\right)}>0.
\end{align*}

Pick any $x,y\in J$. Suppose $x\ne y$. Define $\omega$ to be the longest common word after the fixed prefix $V^+_J$; that is, there exist $m\geq k_{\lambda,f,M,\alpha}+1$, words $U_j\in\{1,\ldots,B\}^{g_j}$ for $j\in\{k_{\lambda,f,M,\alpha}+1,\ldots,m-1\}$, a word $U_m\in\{1,\ldots,B\}^{g_m}$ and a prefix $W$ of $U_m^+$ such that
\begin{align*}
    \omega=U^+_{k_{\lambda,f,M,\alpha}+1}\cdots U^+_{m-1}W
\end{align*}
(the case $m=k_{\lambda,f,M,\alpha}+1$ and $\omega=W$ is allowed). Note that
\begin{align*}
    x=T_{V^+_J\omega}(\xi_x), \qquad
    y=T_{V^+_J\omega}(\xi_y),
\end{align*}
and $a_1(\xi_x)\ne a_1(\xi_y)$. Since $\mathcal{K}_{\lambda,f,M,B}\subset\mathcal{F}_B$, one obtains $a_1(\xi_x),a_1(\xi_y)\in\{1,\dots,B\}$ and $\xi_x,\xi_y\in\mathcal{F}_B\subset[1/\tau_B,1/\varphi]$. Hence, $\xi_x$ and $\xi_y$ belong to two distinct intervals among $\mathscr{J}_B$, so
\begin{align*}
    |\xi_x-\xi_y|\geq \Delta_B.
\end{align*}

By Lemma~\ref{lem:prefix-map-derivative} and the mean value theorem, one obtains
\begin{align}
\label{eq:x-y-lower-local-holder-revised}
    \left|x-y\right|
    =
    \left|T_{V^+_J\omega}(\xi_x)-T_{V^+_J\omega}(\xi_y)\right|
    \ge
    \frac{\Delta_B}{4K(V^+_J\omega)^2}.
\end{align}
Let $\omega^-$ be the word obtained from $\omega$ by deleting all
inserted symbols stage by stage; that is,
\begin{align*}
    \omega^-=U_{k_{\lambda,f,M,\alpha}+1}\cdots U_{m-1}W^-.
\end{align*}
By the definition of $\pi_{\lambda,f,M,B}$ and $W^-$, deleting the inserted symbols from the common prefix $V_J^+\omega$ yields exactly $V_J\omega^-$. Then $\pi_{\lambda,f,M,B}(x)$ and $\pi_{\lambda,f,M,B}(y)$ agree on the prefix $V_J\omega^-$. By the standard cylinder diameter estimate, one obtains 
\begin{align}
\label{eq:pi-x-y-upper-local-holder-revised}
    \left|\pi_{\lambda,f,M,B}(x)-\pi_{\lambda,f,M,B}(y)\right|
    \leq \operatorname{diam}{\mathcal{I}\!\left(V_J\omega^-\right)}
    \leq \frac{1}{K(V_J\omega^-)^2}.
\end{align}

By repeated application of Lemma~\ref{lem:continuant-concatenation}, one obtains
\begin{align*}
    K(V^+_J\omega)
    \leq
    2^{m-k_{\lambda,f,M,\alpha}}
    K(V^+_J)
    \prod_{j=k_{\lambda,f,M,\alpha}+1}^{m-1}K(U^+_j)\,K(W),
\end{align*}
where the empty product is understood to be equal to $1$, and
\begin{align*}
    K(V_J\omega^-)
    \geq
    K(V_J)
    \prod_{j=k_{\lambda,f,M,\alpha}+1}^{m-1}K(U_j)\,K(W^-).
\end{align*}
By applying Lemma~\ref{lem:full-stage-comparison} to the full stages and Lemma~\ref{lem:stage-prefix-comparison} to $W$ respectively, one obtains
\begin{align*}
    K(V^+_J\omega)^\alpha
    &\leq
    2^{(m-k_{\lambda,f,M,\alpha})\alpha}
    K(V^+_J)^\alpha
    \prod_{j=k_{\lambda,f,M,\alpha}+1}^{m-1}K(U^+_j)^\alpha
    K(W)^\alpha \\
    &\leq
    2^{(m-k_{\lambda,f,M,\alpha})\alpha}
    K(V^+_J)^\alpha
    \left(\prod_{j=k_{\lambda,f,M,\alpha}+1}^{m-1}2^{-2\alpha}K(U_j)\right)
    C_{\lambda,f,M,\alpha,2}K(W^-).
\end{align*}
One obtains from $2^{(m-k_{\lambda,f,M,\alpha})\alpha}\prod_{j=k_{\lambda,f,M,\alpha}+1}^{m-1}2^{-2\alpha}=2^{(2+k_{\lambda,f,M,\alpha}-m)\alpha}\leq 2^{2\alpha}$ and the definition of $C_{\lambda,f,M,B,\alpha,0}$ that
\begin{align*}
    K(V^+_J\omega)^\alpha
    \leq
    \frac{2^{2\alpha}C_{\lambda,f,M,\alpha,2}K(V^+_J)^\alpha}{K(V_J)}
    K(V_J\omega^-)
    \leq
    C_{\lambda,f,M,B,\alpha,0}K(V_J\omega^-).
\end{align*}
By combining the above with~\eqref{eq:x-y-lower-local-holder-revised} and~\eqref{eq:pi-x-y-upper-local-holder-revised}, one obtains
\begin{align*}
    \left|\pi_{\lambda,f,M,B}(x)-\pi_{\lambda,f,M,B}(y)\right|
    &\leq \frac{1}{K(V_J\omega^-)^2} 
    \leq \frac{(C_{\lambda,f,M,B,\alpha,0})^2}{K(V^+_J\omega)^{2\alpha}}
    \leq (C_{\lambda,f,M,B,\alpha,0})^2\left(\frac{4}{\Delta_B}\right)^\alpha \left|x-y\right|^\alpha \\
    &\leq C_{\lambda,f,M,B,\alpha}\left|x-y\right|^\alpha,
\end{align*}
where $C_{\lambda,f,M,B,\alpha}\coloneqq (C_{\lambda,f,M,B,\alpha,0})^2({4}/{\Delta_B})^\alpha$.
\end{proof}

\section{Dimension estimate}

\begin{prop}
\label{prop:dimension-lower-final}
Let $\lambda\in\mathbb{N}$ and $f:\mathbb{N}\to\mathbb{N}$ be a function. Suppose $f$ is admissible. Then for any $M,B\in\mathbb{N}$ and $0<\alpha<M/(M+1)$, if $B\geq \lambda+2$ then
\begin{align*}
    \dim{\mathcal{K}_{\lambda,f,M,B}}
    \geq \alpha\dim{\mathcal{F}_B}.
\end{align*}
\end{prop}
\begin{proof}
Let $k_{\lambda,f,M,\alpha}$ be given by Proposition~\ref{prop:local-holder-final}. Since the family $\mathscr{C}_{k_{\lambda,f,M,\alpha}}(\mathcal{K}_{\lambda,f,M,B})$ is finite, there exist $N\in\mathbb{N}$ and $J_1,\dots,J_N\in\mathscr{C}_{k_{\lambda,f,M,\alpha}}(\mathcal{K}_{\lambda,f,M,B})$ such that $\mathscr{C}_{k_{\lambda,f,M,\alpha}}(\mathcal{K}_{\lambda,f,M,B})=\{J_1,\dots,J_N\}$. Note that
\begin{align*}
    \mathcal{K}_{\lambda,f,M,B}=\bigcup_{\nu=1}^N J_\nu.
\end{align*}
Pick any $\nu\in\{1,\ldots,N\}$. One obtains from Proposition~\ref{prop:local-holder-final} that the map
\begin{align*}
    \pi_{\lambda,f,M,B}|_{J_\nu}:
    J_\nu\to \mathcal{I}\!\left(V_\nu\right)\cap \mathcal{F}_B,
\end{align*}
for some finite word $V_\nu\in\bigcup_{n\in\mathbb{N}}\{1,\ldots,B\}^n$, is surjective and $\alpha$-H\"older. By Lemmas~\ref{lem:bounded-prefix-dimension} and~\ref{lem:holder-dimension-estimate}, one obtains
\begin{align*}
    \dim{\mathcal{F}_B}
    =\dim{\left(\mathcal{I}\!\left(V_\nu\right)\cap \mathcal{F}_B\right)}
    \leq\frac{1}{\alpha}\dim{J_\nu}.
\end{align*}
By the finite stability of Hausdorff dimension, 
\begin{align*}
    \dim{\mathcal{K}_{\lambda,f,M,B}}
    = \max_{\nu\in\{1,\ldots,N\}}\dim{J_\nu}
    \geq \alpha\dim{\mathcal{F}_B}.
\end{align*}
\end{proof}

\begin{prop}
\label{prop:eventual-full-dimension}
Let $\lambda\in\mathbb N$ and $f:\mathbb{N}\to\mathbb{N}$ be a function. Suppose $f$ is admissible. Then $\dim{\mathfrak{F}_{\lambda,f}}=1$.
\end{prop}
\begin{proof}
By Propositions~\ref{prop:run-control-final} and~\ref{prop:dimension-lower-final}, for any $M,B\in\mathbb{N}$, if $B\geq \lambda+2$ then
\begin{align*}
    \dim{\mathfrak{F}_{\lambda,f}}
    \geq \dim{\mathcal{K}_{\lambda,f,M,B}}
    \geq\sup_{0<\alpha<M/(M+1)}\alpha\dim{\mathcal{F}_B}
    =\frac{M}{M+1}\dim{\mathcal{F}_B}.
\end{align*}
By the classical fact~\cite{hensley1992continued,jarnik1928metrischen} that $\lim_{B\to\infty}\dim{\mathcal{F}_B}=1$, one obtains that for any $M\in\mathbb{N}$,
\begin{align*}
    \dim{\mathfrak{F}_{\lambda,f}}\geq\frac{M}{M+1}.
\end{align*}
Therefore, $\dim{\mathfrak{F}_{\lambda,f}}\geq1$, and the reverse inequality is trivial.
\end{proof}


\begin{proof}[Proof of Theorem~\ref{thm:main-sn-final}]
Note that $\mathfrak{F}_{\lambda,f}=\bigcup_{N\in\mathbb{N}}\mathfrak{F}_{\lambda,f,N}$, where
\begin{align*}
    \mathfrak{F}_{\lambda,f,N}
    &\coloneqq
    \left\{x\in[0,1)\setminus\mathbb{Q}:
    \lim_{n\to\infty}\frac{L_n(x,\lambda)}{f(n)}=1
    \text{ and }\right.\\
    &\qquad\qquad\left.\text{for any $m\in\mathbb{N}$ and $\mu\in\mathbb{N}\setminus\{\lambda\}$, if $m\geq N$ then }\frac{L_m(x,\lambda)}{L_m(x,\mu)}>1\right\}.
\end{align*}
By Proposition~\ref{prop:eventual-full-dimension} and the countable stability of Hausdorff dimension, $\sup_{N\in\mathbb{N}}{\dim{\mathfrak{F}_{\lambda,f,N}}}=1$.

Pick any $N\in\mathbb{N}$, $y\in\mathfrak{F}_{\lambda,f,N}$. Define $x\coloneqq T_{\lambda^N}(y)$. One claims that for any $\mu\in\mathbb{N}\setminus\{\lambda\}$ and $n\in\mathbb{N}$, $L_n(x,\lambda)>L_n(x,\mu)$. In the case of $n<N$, the first $n$ partial quotients of $x$ are all equal to $\lambda$. One obtains $L_n(x,\lambda)=n>L_n(x,\mu)=0$. In the case of $n\geq N$ and $m\coloneqq n-N<N$, the initial block $\lambda^N$ has length larger than the entire tail segment of length $m$. One obtains $L_n(x,\lambda)\geq N>m\geq L_n(x,\mu)$.

In the case of $n\geq N$ and $m\geq N$, one obtains that every $\mu$-run among the first $N+m$ partial quotients of $x$ is contained entirely in the tail corresponding to the first $m$ partial quotients of $y$. Thus, $L_n(x,\mu)=L_m(y,\mu)$. Since $y\in\mathfrak{F}_{\lambda,f,N}$ and $m\geq N$, one obtains $L_m(y,\lambda)>L_m(y,\mu)$. Every $\lambda$-run among the first $m$ partial quotients of $y$ also occurs among the first $N+m$ partial quotients of $x$; hence, $L_m(y,\lambda)\leq L_{N+m}(x,\lambda)=L_n(x,\lambda)$. Hence, $L_n(x,\lambda)\geq L_m(y,\lambda)>L_m(y,\mu)=L_n(x,\mu)$. The claim is proved for all cases.

Note that for any $m,N\in\mathbb{N}$,
\begin{align*}
    L_m(y,\lambda)
    \leq L_{N+m}(x,\lambda)
    \leq N+L_m(y,\lambda).
\end{align*}
Since $f$ is admissible, for any $N\in\mathbb{N}$ and sufficiently large $m\in\mathbb{N}$, $N\leq f(m)$ and 
\begin{align*}
    1\leq\frac{f(m+N)}{f(m)}\leq\frac{f(m+f(m))}{f(m)}.
\end{align*}
Thus, $\lim_{m\to\infty}f(m+N)/f(m)=1$. By $\lim_{m\to\infty}{L_m(y,\lambda)}/{f(m)}=1$, one obtains from the squeeze theorem that $\lim_{m\to\infty}{L_{N+m}(x,\lambda)}/{f(N+m)}=1$; hence $x\in{\mathfrak E}_{\lambda,f}$, and for any $N\in\mathbb{N}$,
\begin{align*}
    T_{\lambda^N}\!\left(\mathfrak{F}_{\lambda,f,N}\right)
    \subset{\mathfrak E}_{\lambda,f}.
\end{align*}
By Lemma~\ref{lem:prefix-map-derivative}, one obtains that for any $N\in\mathbb{N}$, 
\begin{align*}
    \dim{{\mathfrak E}_{\lambda,f}}
    \geq \dim{T_{\lambda^N}\!\left(\mathfrak{F}_{\lambda,f,N}\right)}
    = \dim{\mathfrak{F}_{\lambda,f,N}}.
\end{align*}
Therefore, $\dim{\mathfrak{E}_{\lambda,f}}\geq1$, and the reverse inequality is trivial.
\end{proof}


\begin{thebibliography}{99}
\bibitem[]{WangWu2011} Wang, Bao-Wei; Wu, Jun. On the maximal run-length function in continued fractions. Annales Univ. Sci. Budapest. Sect. Comput. 34 (2011) 247--268
\bibitem[]{hensley1992continued} Hensley, Doug. Continued fraction {Cantor} sets, {Hausdorff} dimension, and functional analysis. Journal of Number Theory 40 (1992) 336--358
\bibitem[]{jarnik1928metrischen} Jarn{\'i}k, Vojt{\v e}ch. Zur metrischen Theorie der diophantischen Approximationen. Prace Matematyczno-Fizyczne 36 (1928--1929) 91--106
\end{thebibliography}
\end{document}
